\section{Архитектура}
\newcommand{\ArchitectureComponentDiagram}{
\begin{figure}[ht]
\begin{center}
\resizebox{\textwidth}{!}{%
\begin{tikzpicture}[
    component/.style={draw, rounded corners, align=center, minimum width=2.9cm, minimum height=0.9cm, fill=gray!8},
    store/.style={draw, rounded corners, align=center, minimum width=3.1cm, minimum height=0.9cm, fill=blue!7},
    planner/.style={draw, rounded corners, align=center, minimum width=3.3cm, minimum height=0.9cm, fill=green!8},
    dataedge/.style={-{Latex[length=2mm]}, thick},
    controledge/.style={-{Latex[length=2mm]}, dashed, thick},
    backgroundedge/.style={-{Latex[length=2mm]}, dotted, thick}
]
    \node[component] (api) at (0,7.0) {API\\\texttt{read}, \texttt{append}};
    \node[component] (reader) at (-6.2,5.45) {Reader};
    \node[store] (metadata) at (0,5.45) {Metadata\\objects, extents, schemes};
    \node[component] (appender) at (6.2,5.45) {Appender};
    \node[component] (degraded) at (-6.2,3.55) {Degraded\\Read Engine};
    \node[store] (datanodes) at (0,3.55) {Data Nodes\\chunks};
    \node[component] (allocator) at (6.2,3.55) {Allocator\\space, slots};
    \node[planner] (global) at (-4.0,0.65) {GlobalPlanner\\metrics routing\\admission};
    \node[planner] (local) at (1.2,0.65) {LocalPlanner\\temperature, candidate\\in-flight};
    \node[component] (recoder) at (6.2,0.65) {Recoder\\execute transition};

    \draw[dataedge] (api) -- (reader);
    \draw[dataedge] (api) -- (appender);
    \draw[controledge] (reader) -- (metadata);
    \draw[controledge] (appender) -- (metadata);
    \draw[dataedge] (reader) -- (datanodes);
    \draw[dataedge] (reader) -- (degraded);
    \draw[dataedge] (degraded) -- (datanodes);
    \draw[dataedge] (appender) -- (datanodes);
    \draw[controledge] (appender) -- (allocator);
    \draw[backgroundedge] (reader.south west) -- ++(-1.0,-1.05)
        |- node[pos=0.35, left, font=\scriptsize, align=center, fill=white, inner sep=1pt]
        {счётчики чтения\\маршрутизация температуры} (global.west);
    \draw[backgroundedge] (appender.east) -- ++(1.25,0)
        -- node[pos=0.45, right, font=\scriptsize, fill=white, inner sep=1pt] {метрики дозаписи} ++(0,-6.0)
        -| (global.south);
    \draw[backgroundedge] (global.east) to[out=35,in=145] (local.west);
    \draw[backgroundedge] (local.west) to[out=-145,in=-35] (global.east);
    \draw[backgroundedge] (local) -- (recoder);
    \draw[dataedge] (recoder.north west) -- (datanodes.south east);
    \draw[controledge] (recoder.north) -- (allocator.south);

    \node[draw, align=left, rounded corners, fill=white] (legend) at (-6.2,-1.55) {
        \tikz{\draw[dataedge] (0,0) -- (0.6,0);} пользовательские данные\\
        \tikz{\draw[controledge] (0,0) -- (0.6,0);} метаданные и размещение\\
        \tikz{\draw[backgroundedge] (0,0) -- (0.6,0);} фоновое планирование
    };
\end{tikzpicture}%
}
\caption{\label{fig:architecture-components}
Компоненты системы и основные взаимодействия между ними.}
\end{center}
\end{figure}
}

\newcommand{\SlotPlacementArtifacts}{
\begin{table}[ht]
\begin{center}
\caption{\label{tab:slot-placement-example}
Пример слотов размещения для двух объектов.}
\resizebox{\textwidth}{!}{%
\begin{tabular}{|p{0.12\textwidth}|p{0.22\textwidth}|p{0.22\textwidth}|p{0.15\textwidth}|p{0.18\textwidth}|}
\hline
Слот & Узлы данных объекта 1 & Узлы данных объекта 2 & Маска слота & Узлы локализованной избыточности \\
\hline
слот 1 & узлы 1--6 & узлы 1, 3, 6--9 & 0001 & 25--27 \\
\hline
слот 2 & узлы 7--12 & узлы 2, 5, 10, 13--15 & 0010 & 25--27 \\
\hline
слот 3 & узлы 13--18 & узлы 4, 11, 12, 16--18 & 0100 & 25--27 \\
\hline
слот 4 & узлы 19--24 & узлы 19--24 & 1000 & 25--27 \\
\hline
\end{tabular}%
}
\end{center}
\end{table}

\begin{figure}[ht]
\begin{center}
\resizebox{0.82\textwidth}{!}{%
\begin{tikzpicture}[
    slot/.style={draw, rounded corners, align=center, minimum width=2.4cm, minimum height=0.8cm, fill=gray!8},
    edge/.style={-{Latex[length=2mm]}, thick}
]
    \node[slot] (root) at (0,3.2) {1111\\наибольшая полоса};
    \node[slot] (left) at (-3.3,1.7) {0011\\соседние слоты 1--2};
    \node[slot] (right) at (3.3,1.7) {1100\\соседние слоты 3--4};
    \node[slot] (s1) at (-4.8,0) {0001\\слот 1};
    \node[slot] (s2) at (-1.8,0) {0010\\слот 2};
    \node[slot] (s3) at (1.8,0) {0100\\слот 3};
    \node[slot] (s4) at (4.8,0) {1000\\слот 4};

    \draw[edge] (root) -- (left);
    \draw[edge] (root) -- (right);
    \draw[edge] (left) -- (s1);
    \draw[edge] (left) -- (s2);
    \draw[edge] (right) -- (s3);
    \draw[edge] (right) -- (s4);
\end{tikzpicture}%
}
\caption{\label{fig:slot-tree}
Дерево слотов, ограничивающее допустимые объединения и разделения схем.}
\end{center}
\end{figure}

Табл.~\ref{tab:slot-placement-example} и рис.~\ref{fig:slot-tree} показывают
идею слотов размещения. Слоты задаются отдельно для каждого объекта: один и тот
же номер слота у разных объектов может соответствовать разным наборам узлов.
Они нужны не только для уменьшения пространства поиска, но и для того, чтобы
при объединении схем не накапливать широкие полосы, пересекающиеся по одним и
тем же узлам данных. Объединение схем разрешается только для масок, совместимых
по дереву соседних блоков, то есть для соседних поддеревьев одного уровня,
например \(0001\) и \(0010\) или \(0011\) и \(1100\). У таких пар размещение
пользовательских блоков контролируемо непересекающееся, поэтому объединение не
создаёт новую схему с избыточной концентрацией блоков на одних доменах отказа.
}

\newcommand{\TransitionGraphArtifact}{
\begin{figure}[ht]
\begin{center}
\resizebox{\textwidth}{!}{%
\begin{tikzpicture}[
    state/.style={draw, rounded corners, align=center, minimum width=2.2cm, minimum height=0.9cm, fill=gray!8, font=\small},
    generic/.style={draw, rounded corners, align=center, minimum width=2.35cm, minimum height=0.9cm, fill=blue!7, font=\small},
    swap/.style={<->, >=Latex, thick},
    transform/.style={<->, >=Latex, dashed, thick},
    note/.style={align=center, font=\small}
]
    \node[state] (a) at (-5.4,2.4) {\(Hy(3,d,0,0)\)\\реплики};
    \node[state] (b) at (-1.8,2.4) {\(Hy(2,d,1,0)\)\\локальный проверочный блок};
    \node[state] (c) at (1.8,2.4) {\(Hy(1,d,1,1)\)\\LRCC};
    \node[state] (d) at (5.4,2.4) {\(Hy(0,d,1,2)\)\\LRCC};

    \draw[swap] (a) -- (b);
    \draw[swap] (b) -- (c);
    \draw[swap] (c) -- (d);
    \node[note] at (0,3.45) {реплика \(\leftrightarrow\) проверочный блок};

    \node[generic] (l1) at (-4.1,0) {\(Hy(r,d,2l,g)\)};
    \node[generic] (l2) at (0,0) {\(Hy(r,d,l,g)\)};
    \node[generic] (w) at (4.1,0) {\(Hy(r,2d,2l,g)\)};

    \draw[transform] (l1) -- (l2);
    \draw[transform] (l2) -- (w);
    \node[note] at (-2.15,0.9) {разделение/объединение\\локальных групп};
    \node[note] at (2.15,0.9) {разделение/объединение\\полос};
\end{tikzpicture}%
}
\caption{\label{fig:transition-graph}
Семейства допустимых переходов между схемами хранения; запись
\(Hy(r,d,l,g)\) сокращает \(Hy(r,LRCC(d,l,g))\).}
\end{center}
\end{figure}
}

\newcommand{\PlannerArtifacts}{
\begin{figure}[ht]
\begin{center}
\resizebox{\textwidth}{!}{%
\begin{tikzpicture}[
    component/.style={draw, rounded corners, align=center, text width=2.85cm, minimum height=0.9cm, fill=gray!8, font=\small},
    data/.style={draw, rounded corners, align=center, text width=2.85cm, minimum height=0.9cm, fill=blue!7, font=\small},
    datawide/.style={draw, rounded corners, align=center, text width=4.0cm, minimum height=0.9cm, fill=blue!7, font=\small},
    edge/.style={-{Latex[length=2mm]}, thick},
    bidir/.style={<->, >=Latex, thick}
]
    \node[component] (lp) at (-3.6,4.4) {LocalPlanner};
    \node[datawide] (unary) at (-6.1,2.75) {одиночные переходы\\реплика/проверочный блок, изменение \(l\), разделение};
    \node[data] (mergep) at (-0.15,2.75) {пары объединения\\совместимые слоты};
    \node[component] (best) at (-3.25,1.05) {лучший кандидат};
    \node[component] (inflight) at (-6.1,-0.55) {переходы в исполнении};
    \node[component] (recoder) at (-0.15,-0.55) {Recoder};
    \node[component] (gp) at (3.7,4.4) {GlobalPlanner};
    \node[data] (virtual) at (3.7,2.65) {виртуальное состояние\\после допуска};

    \draw[edge] (lp.south west) -- (unary.north);
    \draw[edge] (lp.south east) -- (mergep.north west);
    \draw[edge] (unary.south) -- (best.west);
    \draw[edge] (mergep.south) -- (best.east);
    \draw[edge] (lp.east) to[out=20,in=160]
        node[above, font=\scriptsize, fill=white, inner sep=1pt] {кандидат} (gp.west);
    \draw[edge] (gp.west) to[out=-160,in=-20]
        node[below, font=\scriptsize, fill=white, inner sep=1pt] {допуск} (lp.east);
    \draw[bidir] (gp) -- (virtual);
    \draw[edge] (best) -- node[left, font=\scriptsize, fill=white, inner sep=1pt] {учёт} (inflight);
    \draw[edge] (best) -- node[right, font=\scriptsize, fill=white, inner sep=1pt] {запуск} (recoder);
\end{tikzpicture}%
}
\caption{\label{fig:planner-scheme}
Схема взаимодействия локальных и глобального планировщиков.}
\end{center}
\end{figure}
}


В разделе~\ref{sec:theory-hybrid} была введена модель
\(Hy(r, LRCC(d,l,g))\), в которой реплики, локальные проверочные блоки и
глобальные проверочные блоки рассматриваются как части одного пространства
избыточности. В этой главе описывается архитектура системы, использующей это
пространство состояний для фонового изменения схемы хранения уже записанных
данных. Акцент делается на границах компонентов, инвариантах размещения,
допустимых переходах и принципе работы планировщиков.

\subsection{Архитектурная постановка}

Система рассматривается в постановке с дозаписью: после записи пользовательские
данные не изменяются на месте, но способ их избыточного хранения может
изменяться в фоне. Поэтому в архитектуре разделяются два контура.
Пользовательский контур обслуживает чтение логических диапазонов объекта и
должен давать корректный ответ в текущем состоянии схемы. Фоновый контур
изменяет состояние схемы, выбирая между репликами, локальными и глобальными
проверочными блоками так, чтобы улучшать баланс между занятым местом и
ожидаемой стоимостью чтения.

Из этого следует разделение ответственности между компонентами:
\begin{itemize}
    \item интерфейс пользовательских операций работает с объектами и
    логическими диапазонами;
    \item слой метаданных задаёт согласованное отображение логических диапазонов
    в текущие схемы хранения и физические блоки;
    \item перекодировщик исполняет выбранный переход между схемами;
    \item локальные и глобальный планировщики выбирают полезные переходы на
    основе метрик кластера.
\end{itemize}

Такое разделение выбрано вместо монолитного компонента, который одновременно
обслуживал бы чтения, принимал решения о перекодировании и менял размещение.
Причина в разных требованиях к этим операциям. Чтение чувствительно к задержке
и должно использовать уже зафиксированную схему. Перекодирование может
выполняться дольше, но обязано сохранять требуемую отказоустойчивость во время
перехода. Планирование требует агрегированного состояния кластера и не должно
встраиваться в критический путь пользовательского чтения.

\ArchitectureComponentDiagram

\subsection{Размещение и совместимость переходов}

Размещение должно учитывать не только текущую отказоустойчивость, но и будущую
стоимость переходов. Основной инвариант размещения состоит в \emph{разнесении
по доменам отказа}: блоки одной схемы, совместно необходимые для восстановления,
не должны попадать в один домен отказа. В данной работе доменом отказа считается
узел кластера. При выполнении этого условия потеря одного узла не уничтожает
сразу несколько критичных блоков одной схемы.

Для проверочных блоков выбирается объектно-локализованное размещение.
Альтернативой было бы равномерно распределять все блоки по кластеру. Такой
подход лучше симметризует нагрузку фонового восстановления при отказе узла, но
делает будущие переходы дороже: совместимые проверочные блоки одного объекта
часто оказываются на разных узлах. Объектная локализация, наоборот, повышает
вероятность дешёвых переходов, особенно для операций над проверочными блоками.
Цена этого решения --- риск всплеска восстановления при отказе узла, на котором
накопилось много проверочных блоков одного или нескольких объектов. Поэтому
локализация должна сочетаться с разнесением блоков схемы по доменам отказа,
учётом свободного места и мониторингом потенциальной цены восстановления.

Для операций объединения и разделения схем размещение должно быть проверяемым
на совместимость. Объект получает набор слотов размещения: логических групп
доменов отказа, которые используются для проверки, можно ли объединить две
схемы без перемещения пользовательских данных. Чтобы ограничить число
допустимых пар, используется правило дерева соседних блоков: объединять можно
только схемы, чьи множества слотов образуют допустимые соседние поддеревья. Это
не универсальное свойство кодов, а выбранное ограничение пространства поиска.

Совместимость размещения является необходимым, но не единственным условием.
Схемы также должны быть кодово совместимы. В данной архитектуре это означает,
что схемы принадлежат одному объекту, имеют одинаковую или приводимую к общей
цели сигнатуру \(Hy(r, LRCC(d,l,g))\), а для операций с локальными проверочными
блоками имеют одинаковый размер локальной группы \(d/l\) после возможного
изменения числа локальных групп. Кроме того, обе схемы должны иметь общий
целевой переход: например, объединение двух схем в одну более широкую схему
или обратное разделение одной схемы на две.

\SlotPlacementArtifacts

\subsection{Пространство состояний и переходов}

Архитектура не фиксирует одну цепочку состояний. Состояние схемы задаётся
моделью \(Hy(r, LRCC(d,l,g))\), введённой в разделе~\ref{sec:theory-hybrid}.
В архитектуре эта модель используется как пространство допустимых состояний:
каждый фоновый переход должен переводить допустимую схему в другую допустимую
схему и сохранять инварианты размещения, локальных групп и независимости
глобальных проверочных блоков.

Планировщик выбирает не абстрактное конечное состояние жизненного цикла, а
допустимый переход из текущего состояния. К таким переходам относятся:
\begin{itemize}
    \item обмен одной или нескольких реплик на локальные и глобальные
    проверочные блоки;
    \item обратный обмен проверочных блоков на реплики;
    \item изменение числа локальных групп в два раза при сохранении допустимого
    размера локальной группы;
    \item разделение одной схемы на две схемы;
    \item объединение двух совместимых схем одного объекта.
\end{itemize}

Переходы в сторону большего числа проверочных блоков обычно выгодны при
охлаждении данных или росте давления по месту. Переходы в сторону большего
числа реплик выгодны при росте температуры и ожидаемой стоимости чтения.
Однако это не жёсткая цепочка: планировщик может выбрать прямой или составной
переход между допустимыми состояниями, если его полезность выше с учётом
стоимости фоновой работы.

\TransitionGraphArtifact

\subsection{Фоновое перекодирование}

Перекодировщик является исполнительным компонентом: он получает выбранный
переход и выполняет физическую смену избыточности. Его главный инвариант ---
сохранять требование восстановления после любых трёх отказов на всём протяжении
перехода. Поэтому новая избыточность должна стать доступной до удаления той
избыточности, которую она заменяет, а переключение метаданных должно иметь
атомарную точку фиксации.

Для переходов от реплик к проверочным блокам перекодировщик читает достаточно
блоков данных, вычисляет новые локальные или глобальные проверочные блоки,
размещает их согласно политике размещения и только затем удаляет заменяемые
реплики. Для обратных переходов он создаёт новые реплики из доступной реплики,
а если реплик в схеме не осталось --- из LRCC-блоков данных. Эта асимметрия
важна для оценки стоимости перехода: нагрев схемы после полного ухода в
кодовое представление дороже, чем добавление реплики при наличии уже доступной
реплики.

Если исполнитель прерывается во время перехода, корректность должна
обеспечиваться протоколом фиксации метаданных: до точки фиксации старая схема
остаётся действующей, после точки фиксации действующей считается новая схема.
Журналирование, идемпотентность повторного выполнения и точный способ очистки
частично записанных блоков относятся к реализации этого протокола, но
архитектурный инвариант остаётся тем же: удаление старой избыточности не должно
предшествовать доступности новой.

\subsection{Чтение и восстановление}

Пользовательское чтение выбирает самый дешёвый доступный способ получить
нужный логический блок. Если доступна реплика, чтение использует её. Если
реплик нет или они недоступны, чтение обращается к блоку данных в
LRCC-представлении. Такой порядок соответствует назначению гибридных схем:
реплики уменьшают стоимость чтения, а LRCC-представление обеспечивает
компактность и восстановимость.

Если прямое чтение невозможно, выполняется чтение с восстановлением. Оно
обслуживает текущий пользовательский запрос и не меняет физическое состояние
хранилища. Сначала используется локальный проверочный блок, если недоступный
блок может быть восстановлен внутри локальной группы. Если локального
проверочного блока недостаточно, используются глобальные проверочные
уравнения. Для допустимых схем из семейства \(Hy(r, LRCC(d,l,g))\)
восстановление возможно до целевого уровня отказов при условии, что остаётся
достаточно независимых уравнений и размещение по-прежнему обеспечивает
разнесение совместно необходимых блоков по разным доменам отказа.

Фоновое восстановление физически утраченного блока отделяется от чтения с
восстановлением. Оно относится к эксплуатации кластера: нужно обнаружить
длительно недоступный блок, выбрать новое размещение, записать восстановленный
блок и обновить метаданные. В данной работе этот процесс не является переходом
в пространстве схем и не входит в задачу выбора формы избыточности.

\subsection{Метрики состояния кластера}

Планировщик использует две основные метрики: давление по занятому месту и
давление по ожидаемой стоимости чтения. Температура хранится на уровне схемы,
поскольку именно схема является единицей планирования. Температура на уровне
отдельных логических блоков может точнее описывать горячие участки больших
объектов, но она усложняет агрегацию при объединении и разделении схем.
Поэтому для данной архитектуры используется температура уровня схемы, а более
детальное представление может быть добавлено как уточнение этой метрики.

Обновление температуры выполняется асинхронно. Путь чтения накапливает счётчики
обращений, ответственный локальный планировщик получает эти счётчики и
обновляет экспоненциально взвешенное скользящее среднее (EWMA). Между
обновлениями применяется ленивое затухание: значение пересчитывается при
обращении к схеме или при планировочном проходе.

Вероятности отказов оцениваются отдельно от температуры. Величина \(\hat q\)
означает сглаженную долю недоступных доменов отказа в кластере. Для конкретной
схемы из неё строится \(p_f\) --- вероятность того, что среди доменов отказа,
занятых этой схемой, недоступны ровно \(f\) доменов. Эти вероятности входят в
ожидаемую стоимость чтения \(E[io]\), определённую в
разделе~\ref{sec:theory-read-cost}. Итоговая полезность перехода вычисляется
через скалярную функцию давления из раздела~\ref{sec:theory-temperature}, то
есть через \(P_{space}\), \(P_{io}\) и \(Q\).

\subsection{Локальное и глобальное планирование}

Локальный планировщик владеет объектно-локальной группой: группой записей схем
одного объекта, сохраняющей локальность для поиска совместимых переходов. Такое
владение нужно прежде всего для операций объединения: совместимые схемы почти
всегда должны сравниваться внутри одного объекта, а не по всему кластеру. Один
планировщик может владеть несколькими объектно-локальными группами, но каждая
схема в каждый момент времени имеет одного ответственного планировщика.

Во время планировочного прохода локальный планировщик получает текущие метрики
кластера, обновляет температуры принадлежащих схем с учётом ленивого затухания
и строит кандидаты переходов. Для одиночных переходов он оценивает полезность
каждого допустимого изменения состояния. Для объединения он группирует схемы по
профилю совместимости и рассматривает только пары, удовлетворяющие ограничениям
размещения и кодовой совместимости. Такой поиск остаётся локальным для объекта
и не превращается в полный попарный перебор всех схем кластера. В конце прохода
локальный планировщик сохраняет один лучший кандидат, который может быть
одиночным переходом или точной парой объединения.

Глобальный планировщик работает в режиме опроса. В начале итерации он опрашивает
локальные планировщики и получает от каждого не более одного лучшего кандидата.
Если схема уже занята перекодированием, ответственный локальный планировщик сам
не возвращает для неё кандидата. Полученные кандидаты сортируются по полезности;
такая сортировка допустима, поскольку число ответственных планировщиков
выбирается порядка числа узлов или разделов планирования, а не порядка числа
всех схем. Далее глобальный планировщик последовательно выдаёт разрешения на
выполнение и не блокируется на завершении физического перекодирования.

Чтобы после первого допущенного перехода не ждать его фактического завершения,
глобальный планировщик ведёт виртуальное состояние кластера. После каждого
допуска он применяет ожидаемый эффект перехода к виртуальным метрикам и
оценивает следующие кандидаты так, как если бы все уже допущенные переходы
повлияли на кластер. Это позволяет запускать несколько переходов подряд, но не
выдавать решения, конфликтующие по влиянию на оценку полезности. Учёт уже
допущенных, но ещё не завершённых переходов ведётся у локального планировщика
как состояние переходов в исполнении. Выдача разрешений останавливается, когда
исчерпаны полезные кандидаты или достигнут предел одновременно выполняемых
фоновых переходов.

Такой арбитраж не полностью исключает осцилляции, но снижает их вероятность.
Если множество локальных планировщиков одновременно действовало бы по одному
устаревшему снимку, кластер мог бы слишком резко уйти в одну сторону, например
в чрезмерно компактные схемы, а затем начать выполнять обратные переходы из-за
роста давления по чтению. Последовательное применение переходов к виртуальному
состоянию позволяет остановиться до такого переизбыточного эффекта.
Алгоритмический вид планировочных проходов для реализованного прототипа
приведён в разд.~\ref{sec:implementation-planners}.

\PlannerArtifacts

\subsection{Границы архитектурной модели}

Архитектура опирается на несколько ограничений постановки. Во-первых, данные
считаются неизменяемыми после записи, поэтому пользовательские обновления и
перезапись диапазонов не рассматриваются. Во-вторых, фоновое восстановление
физически утраченных блоков отделено от адаптивного перекодирования и не
является переходом между схемами. В-третьих, функция выбора перехода остаётся
скалярной: она учитывает занятое место, ожидаемую стоимость чтения и стоимость
фоновой работы. Число операций диска, сетевой трафик, число задействованных
узлов и процессорная работа могут использоваться как диагностические метрики и
как основа модели задержки, но не образуют отдельную взвешенную функцию
полезности.

\pagebreak
